\begin{solapartat}
\punts{0,7}
\figura[0.4]{sol_fig1.png}

Si els nodes no estan ben identificats es restarà 0,2~p.

Si els ventres no estan ben identificats es restarà 0,2~p.

Si no s'identifica la primera ressonància com el primer harmònic es restarà 0,1~p.

Si no s'identifica la segona ressonància com el tercer harmònic es restarà 0,2~p.

No es restarà el fet que no escalin correctament la longitud L.

\punts{0,15} Podem comprovar que com que és un tub obert per un costat a la part superior tenim un ventre, per tant pel primer harmònic tenim:
\[ L = \frac{\lambda}{4} \;\Rightarrow\; \lambda_1 = 4L = 0,7720\ \mathrm{m} \]
Pel tercer harmònic tenim
\[ L = \frac{\lambda_3}{4} + \frac{\lambda_3}{2} = \frac{3\lambda_3}{4} \;\Rightarrow\; \lambda_3 = \frac{4}{3}L = 0,7733\ \mathrm{m} \]

\punts{0,2} Podem comprovar que en els dos casos obtenim la mateixa longitud d'ona, la raó és que en ambdós casos la freqüència de l'ona ressonant és la mateixa, 442~Hz, la velocitat de propagació és la mateixa, $c$, llavors com la longitud d'ona també ha de ser la mateixa:
\[ \lambda = \frac{c}{f} \]

\punts{0,2}
\figura[0.22]{sol_fig2.png}

El cinquè harmònic és l'indicat a la figura
\[ L = \frac{\lambda_5}{4} + \lambda_5 = \frac{5\lambda_5}{4} = \frac{5\cdot 0,7727}{4} = 0,966\ \mathrm{m} \]
Com a valor de longitud d'ona hem agafat el valor mitjà del primer i tercer harmònics, però es pot considerar vàlida la resposta amb el valor de $\lambda_1$ o $\lambda_3$.
\end{solapartat}

\begin{solapartat}
\punts{0,65} La velocitat del so és:
\[ c = \lambda f = 0,7727\cdot 442 = 340\ \frac{\mathrm{m}}{\mathrm{s}} \]

\punts{0,6} No variarà perquè l'ona que percebem és la que es propaga a través de l'aire, i per tant la velocitat de propagació en tots dos casos serà la mateixa, 340~m/s.
\end{solapartat}
